\section{Planned mechanism comparison}

\wip{No mechanism has been fitted and no mechanism result is claimed in this
draft. This section states the comparison the completed paper must make.}

The mechanism layer separates two questions that are easily conflated:
\begin{enumerate}
  \item \textbf{Composition:} which fertility-related dispositions or social
  types become more common across generations?
  \item \textbf{Environment:} how many births does a person of a given type
  realize under changing economic, institutional, and cultural conditions?
\end{enumerate}

Intergenerational persistence can shift composition while development,
opportunity cost, partnership patterns, and institutions shift the mapping
from disposition to realized fertility. A completed model must therefore allow
selection and environment to move simultaneously and in opposite directions.
It must also distinguish genetic transmission, family or cultural transmission,
group retention, intermarriage, and assimilation rather than compress those
channels into one uninterpretable coefficient.

Every mechanism parameter will be accompanied by a provenance table stating
its independent evidence source, the outcome used to estimate it, its prior,
and whether it remains a scenario choice. A parameter may not be fitted to the
same outcome that it is later said to explain. Period total fertility may be
reported descriptively, but it will not be used as the scoring target because
changes in birth timing can distort it. Completed-cohort fertility and parity
progression are needed to distinguish tempo from quantum
\citep{bongaarts1998}.
